\chapter{实例效果分析}


\section{学生反馈}

本实例在福州市某高中高一年段的两个班级进行应用，针对本课堂的学生进行课堂评价、教学效果等多方面的问卷调查，通过线下进行问卷收集。问卷包含三个部分，第一个部分为被调查学生的基本信息，包括性别与班级；第二部分是问卷的主体部分，对被调查学生的教学效果与目标达成进行了调查，采用李克特式五点记分制，共10道题；第三部分为探究性学习课堂偏好调查，共2道题。


\subsection{被调查学生变量频率分析}

\begin{table}[htbp]
  \centering
  \caption{被调查学生变量频率分析}
  \label{被调查学生变量频率分析}
  \begin{tblr}{
    width = \textwidth,
    hline{1,2,Z} = {solid},
    colspec = {X[1,c]X[2,c]X[1,c]X[1.5,c]X[1.5,c]X[1.5,c]},
    cell{2}{1} = {r=2}{m},
    cell{4}{1} = {r=2}{m},
    cell{2}{5} = {r=2}{m},
    cell{2}{6} = {r=2}{m},
    cell{4}{5} = {r=2}{m},
    cell{4}{6} = {r=2}{m},
  }
    变量 & 选项    & 频率 & 百分比  & 平均值  & 标准差 \\
    性别 & 男     & 43 & 52\% & 1.48 & 0.5 \\
      & 女     & 40 & 48\% &      &     \\
    班级 & 高一11班 & 41 & 49\% & 1.51 & 0.5 \\
      & 高一12班 & 42 & 51\% &      &   
  \end{tblr}
\end{table}


根据表~\ref{被调查学生变量频率分析} 的结果可以看出被调查学生的数值特征，此次问卷调查的对象为福州市某高中高一两个班级的学生，均是本文中教学案例的实验者。本次调查共收到95份问卷，经笔者筛查得到有效问卷83份，回收率87.37\%。


\subsection{问卷信度检险}


问卷量表~\ref{问卷信度分布表} 调查数据输入SPSS27.0进行信度分析。

\begin{table}[htbp]
  \centering
  \caption{问卷信度分布表}
  \label{问卷信度分布表}
  \begin{tblr}{
    width = \textwidth,
    hline{1,2,Z} = {solid},
    colspec = {X[2,c]X[3,c]X[1,c]}
  }
  项目        & 克隆巴赫Alpha系数 & 项目数 \\
  知识与技能层面   & 0.846       & 3   \\
  过程与方法层面   & 0.902       & 4   \\
  情感、态度与价值观 & 0.901       & 3   \\
  问卷量表      & 0.944       & 10 
  \end{tblr}
\end{table}


问卷整体的克隆巴赫Alpha系数为0.944，达到了很好的信度水平，其中各项教学目标量表的克隆巴赫Alpha系数均高于0.8，因此问卷信度良好。


\subsection{描述性统计分析}

\begin{table}[htbp]
  \centering
  \caption{问卷量表的描述性分析}
  \label{问卷量表的描述性分析}
  \begin{tblr}{
    width = 1.1\textwidth,
    hline{1,2,5,9,Z} = {solid},
    colspec = {X[7em,c]X[19em,c]X[1,c]X[1,c]X[1,c]},
    cell{1}{1} = {c=2}{c},
    cell{2}{1} = {r=3}{m},
    cell{5}{1} = {r=4}{m},
    cell{9}{1} = {r=3}{m},
    cell{2}{5} = {r=3}{m},
    cell{5}{5} = {r=4}{m},
    cell{9}{5} = {r=3}{m},
  }
  项目      &                     & 平均值  & 标准差   & 总均值    \\
  知识与技能层面 & 知识点的记忆与掌握           & 3.93 & 0.973 & 3.12   \\
          & 如遇类似习题能够举一反三        & 3    & 1.126 &        \\
          & 能够将该知识点应用在实际生活中     & 2.43 & 1.014 &        \\
  过程与方法层面 & {在老师引导下，能够推广习题得出新的命题} & 3.37 & 0.946 & 3.2275 \\
          & 掌握“发现—猜想—解决”的思路与方法  & 3.58 & 0.813 &        \\
          & {能够与小组成员相互合作，共同解决问题}  & 2.94 & 0.942 &        \\
          & 能够纠正自己或他人的错误命题与方法   & 3.02 & 1     &        \\
  {情感、态度与价值观层面}   & 课堂具有启发性与趣味性         & 4.04 & 0.756 & 3.22   \\
    & 能够主动发现、探究、学习知识      & 2.75 & 0.948 &        \\
          & {能够反思问题的探究过程并进行检查修正}  & 2.87 & 0.947 & 
  \end{tblr}
\end{table}

表~\ref{问卷量表的描述性分析} 说明，过程与方法层面自评分数较高，其中“掌握‘发现—猜想—解决’的思路与方法”评分最高；知识与技能层面均分较低，为3.12，说明该案例在知识讲解与技能掌握方面有所欠缺，但在研究方法与路径方面成果显著。

在知识与技能模块，对于“知识点的记忆与掌握”，有34.9\%和33.7\%的学生选择了“4”和“5”，没有学生选择“1”，说明学生总体反映对知识点的记忆和理解程度较高；对于“如遇类似习题能够举一反三”，约有35\%学生选择“4”，同时也注意到约27\%学生选择“2”，这说明对于举一反三这一技能在本教学案例体现出不足之处，有部分学生不确定自己是否能够对类似的习题进行同样的分析；对于“能够将该知识点应用在实际生活中”，有36.14\%的学生选择“2”，说明多数学生没有将函数与实际问题联系起来，缺乏该方面的链接。

在过程与方法模块，前两个项目均分较高，说明大多数学生能够在教师引导下进行探究性学习，并且掌握探究性学习的一般路径；对于小组合作方面，约有50\%的学生选择“3”，说明小组合作的效果处于中等水平，还需要进一步优化；针对反思纠错方面，学生总体表现出的水平不是太强。

在情感、态度与价值观模块，对于课堂的趣味性与启发性，大多数学生给予了高度评价，探究性学习在习题课中能够激发起学生学习的兴趣，并且给予思维的启示；但是大多数学生在主动进行探究性学习的能力上稍有欠缺。

\begin{table}[htbp]
  \centering
  \caption{习题课类型偏好分析}
  \label{习题课类型偏好分析}
  \begin{tblr}{
    width = \textwidth,
    hline{1,2,Z} = {solid},
    colspec = {X[2,c]X[3,c]X[1,c]}
  }
    选项     & 小计 & 比例      \\
    传统习题课  & 8  & 9.64\%  \\
    探究性习题课 & 24 & 28.92\% \\
    二者结合   & 40 & 48.19\% \\
    并无偏好   & 11 & 13.25\%
  \end{tblr}
\end{table}


在习题课类型偏好方面，从表~\ref{习题课类型偏好分析} 可以看出约有五成学生喜好探究性学习与传统讲解学习相结合的习题课堂，也有不少同学喜爱探究性习题课堂，这说明传统模式的习题课堂已逐渐在学生心中淘汰，大部分学生希望寻求趣味性高、思维性强的习题课。


\begin{table}[htbp]
  \centering
  \caption{探究性习题课堂教师指导偏好分析}
  \label{探究性习题课堂教师指导偏好分析}
  \begin{tblr}{
    width = \textwidth,
    hline{1,2,Z} = {solid},
    colspec = {X[2,c]X[3,c]X[1,c]}
  }
    选项    & 小计 & 比例      \\
    知识的辅导 & 60 & 72.29\% \\
    方法的引导 & 70 & 84.34\% \\
    意识的培养 & 25 & 30.12\% \\
    情感的激励 & 38 & 45.78\% \\
    思维的突破 & 37 & 44.58\% \\
    其他    & 0  & 0\%    
  \end{tblr}
\end{table}

在探究性习题课堂教师指导偏好方面，从表~\ref{探究性习题课堂教师指导偏好分析} 可以看出超过八成的学生期望从探究性习题课上获得数学方法的引导，希望获得知识辅导的学生不在少数，同时也有约50\%的学生期望获得教师在情感方面的鼓励和数学思维的培养与突破。由数据可得，大部分学生偏好习题课堂中的知识与方法的传授，同时该偏好也可能侧面反映了现有的习题课所缺失的教师指导。

除此之外，笔者还通过学生课后反馈得知，在对比不同习题课类型的花费时间上，学生普遍认为在讲解相同内容的情况下，探究性习题课堂所花费的时间比现有的习题课要多。但在课堂节奏紧凑且前期准备充分的情况下，该时间可以改善。


\subsection{差异性分析}

\begin{table}[htbp]
  \centering
  \caption{各维度在班级上的差异性分析}
  \begin{tblr}{
    hline{1,2,Z} = {solid},
    colspec = {X[3,c]X[2,c]X[1,c]X[1,c]X[1,c]},
    cell{2}{1} = {r=2}{m},
    cell{4}{1} = {r=2}{m},
    cell{6}{1} = {r=2}{m},
  }
    变量 & 班级 & 数字 & 平均值 & 标准差 \\
    知识与技能层面 & 高一（11）班 & 41 & 8.93 & 2.944 \\
      & 高一（12）班 & 42 & 9.79 & 2.465 \\
      过程与方法层面 & 高一（11）班 & 41 & 12.71 & 3.052 \\
      & 高一（12）班 & 42 & 13.12 & 3.466 \\
      情感、态度与价值观层面 & 高一（11）班 & 41 & 9.39 & 2.607 \\
      & 高一（12）班 & 42 & 9.9 & 2.239 \\
  \end{tblr}
\end{table}

经笔者了解到，高一（11）班为平行班级，高一（12）班为培优班，由于两个班级的生源水平与教学方向不同，因此笔者考虑对班级与三维目标作差异性检验。由表5-6数据可知，在每个层面，12班所得均分均高于11班，说明学生学业水平的高低会影响课堂目标实现的效果，同时班级的差异性恰好说明该量表的有效性良好。在知识与技能层面，均分差异为0.86，差值最大，说明学生的水平在此目标上的影响最大化；反之，在过程与方法层面差值最小，学生在此目标的反馈较弱，说明教师更应该重视该目标的培养。



\section{教师反馈}

在教学实例实施后，笔者针对有关探究性学习在习题课中应用的看法与建议对旁听教师进行访谈，访谈的结果整理如下：

\begin{enumerate}[1.]
  \item 结合本实例的实际表现，您认为探究性学习在习题课中的应用有哪些优势吗？

  旁听教师普遍认为，在以往的习题课堂上主要采用的是教师讲授的教学形式，以教师为主导地位，学生的主体地位并未凸显，相反，在探究性习题课上学生能够主动地去发现知识，并且证明知识。高一学生对函数的性质这一内容在理解上还比较困难，因此让学生去探索有利于知识的记忆与理解。许多习题都可以采用这种教学形式，挖掘本应该由教师讲述的知识点，真正让学生参与进来有所收获。
  
  \item 对比本实例与以往课堂，您认为学生上课的表现有何变化吗？
  
  旁听教师反应相比于以往课堂，学生在课堂的参与程度和踊跃程度大大提升，课堂气氛相对较活跃，学生能够大胆地进行小组讨论，紧跟教师教学节奏，且气氛热烈。若平时课堂也能够如此，那么教师与学生相互配合的效率将会高很多。许多教师还反应，探究性学习在课堂中的应用有利于增进师生关系，课堂活跃度往往和师生关系十分密切，在实例课堂上，学生的参与度和踊跃度都比平时课堂高，因此对师生关系起到积极作用。
  
  \item 您认为探究性学习在习题课中的应用对学生解题有何帮助？
  
  在习题教学中，探究性学习的应用促使了习题的核心思维的体现，学生能够得到数学思维的培养与创新能力的提升。通过教师对于探究性学习的一般方法的指导，学生能够在面对类似习题、命题和问题时能够运用相同的解题路径，从而提高学生的解题能力。许多教师还反应，当今的多数习题考查的不是单一的知识点，也不是简单地举一反三，而是需要通过一定研究方法得到结论的综合型题目，这也是新兴的高考命题方向，因此探究性学习应该被教师重视，并且运用到日常课堂中来。
\end{enumerate}